\abstract

Работа посвящена верификации соответствия трасс исполнения программы формальной
спецификации на языке TLA+.

Целью работы является разработка методики сопоставления трассы исполнения
программы с формальной спецификацией, формализация соответствующей
математической постановки задачи и обоснование её решения.

В работе предложена методика, основанная на инструментировании программной
реализации, построении трассы наблюдаемых шагов исполнения и её интерпретации
как системы ограничений на переходы формальной модели. Сформулирована
математическая постановка задачи в терминах системы переходов, множества
допустимых поведений модели и множества поведений, согласованных с трассой.
Показано, что рассматриваемая задача является задачей разрешимости. Для её
решения введена последовательность множеств достижимых состояний, согласованных
с префиксами трассы, и на этой основе получен алгоритм проверки согласованности
трассы с моделью.

Практическая реализация предложенного подхода выполнена средствами TLA+ и TLC
на примере спецификации протокола двухфазной фиксации транзакции.

\textbf{Ключевые слова:} TLA+, TLC, формальная верификация, система переходов,
трасса исполнения, распределённые системы.
